\section{The complete transition table}
\label{app:rules}

Our automaton has a compact rule set, unlike the tiling-specific automata, whose appendices list hundreds of
rules because each enumerates every rotation-invariant full-neighbourhood configuration of its tiling. Ours
reads a cell's role, its dynamic symbol, its switch labels and selection, and the symbols of its linked
neighbours, so the whole transition function is the small set below, the same on every tiling. We list it in
full, generated directly from the automaton, so that it is as explicit as those appendices, only shorter.

\subsection{Motion of the locomotive}

On a track or switch cell the head and trailing cell move by Table~\ref{tab:enum-motion}. The head reads which
linked neighbour holds its trailing cell, and pushes the head into the other (forward) linked neighbour.

\begin{table}[ht]
\centering
\begin{tabular}{@{}lll@{}}
\toprule
the cell is & condition & it becomes \\
\midrule
empty & always & empty \\
$A$ (trailing) & always & $C$ \\
$H$ (head) & always & $A$, and pushes $H$ into its forward linked cell \\
$C$ (clear) & a linked neighbour is $H$ and its forward cell is this one & $H$ \\
$C$ (clear) & otherwise & $C$ \\
\bottomrule
\end{tabular}
\caption{The motion rules in full. The forward cell of a head is its linked neighbour that does not hold the
trailing cell, which is unique by Lemma~\ref{lem:direction}.}
\label{tab:enum-motion}
\end{table}

\subsection{The switches, every case}

Table~\ref{tab:enum-switch} is the complete switch transition, every kind, every entry side, and every
selection, generated by running each case on the automaton. The columns are the switch kind, the branch the
head entered from (the trunk $u$, or a branch $a$ or $b$), the current selection, the branch the head exits by,
and the selection after. A passive crossing, entered from a branch, always exits the trunk. An active crossing,
entered from the trunk, exits the selected branch. The flip-flop flips its selection only on the active
crossing. The memory switch sets its selection to the branch of a passive crossing.

\begin{table}[ht]
\centering
\begin{tabular}{@{}lcccc@{}}
\toprule
kind & entered from & selection & exits by & new selection \\
\midrule
fixed & $u$ & $a$ & $a$ & $a$ \\
fixed & $u$ & $b$ & $b$ & $b$ \\
fixed & $a$ & $a$ & $u$ & $a$ \\
fixed & $a$ & $b$ & $u$ & $b$ \\
fixed & $b$ & $a$ & $u$ & $a$ \\
fixed & $b$ & $b$ & $u$ & $b$ \\
\midrule
flip-flop & $u$ & $a$ & $a$ & $b$ \\
flip-flop & $u$ & $b$ & $b$ & $a$ \\
flip-flop & $a$ & $a$ & $u$ & $a$ \\
flip-flop & $a$ & $b$ & $u$ & $b$ \\
flip-flop & $b$ & $a$ & $u$ & $a$ \\
flip-flop & $b$ & $b$ & $u$ & $b$ \\
\midrule
memory & $u$ & $a$ & $a$ & $a$ \\
memory & $u$ & $b$ & $b$ & $b$ \\
memory & $a$ & $a$ & $u$ & $a$ \\
memory & $a$ & $b$ & $u$ & $a$ \\
memory & $b$ & $a$ & $u$ & $b$ \\
memory & $b$ & $b$ & $u$ & $b$ \\
\bottomrule
\end{tabular}
\caption{The complete switch transition, all eighteen cases, generated from the automaton. Note the flip-flop
rows, entry from the trunk flips the selection, entry from a branch leaves it unchanged. Note the memory rows,
entry from branch $a$ leaves or sets the selection to $a$, entry from branch $b$ to $b$.}
\label{tab:enum-switch}
\end{table}

\subsection{The track-building rule}

The one further rule, used only for strong universality, is the builder of Table~\ref{tab:build}, a head at the
end of a register with no clear track ahead turns its deepest empty neighbour into track and advances into it.
With Tables~\ref{tab:enum-motion}, \ref{tab:enum-switch}, and~\ref{tab:build}, the automaton is completely
specified. This is the entire rule set, the same on every regular hyperbolic tiling.
